\Askhsh{24275_SOLUTION}{1}
\begin{ylh}{1}
    \item [Ύλη από εικόνα]
\end{ylh}
\begin{alist}
\item Επειδή
\[ \lim_{x \to +\infty} (f(x) - (-x+1)) = \lim_{x \to +\infty} \left(-x+1+\frac{1}{e^x} + x-1\right) = \lim_{x \to +\infty} \frac{1}{e^x} = 0 \]
έπεται ότι η ευθεία $y=-x+1$ είναι πλάγια ασύμπτωτη της γραφικής παράστασης της $f$ στο $+\infty$.

\item Είναι, για κάθε $x \in \mathbb{R}$,
\[ f'(x) = \left(-x+1+\frac{1}{e^x}\right)^{\prime} = (-x)^{\prime}+(1)^{\prime}+(e^{-x})^{\prime} = -1+0+e^{-x}(-x)^{\prime} = -1-e^{-x}. \]
Άρα $f'(x)<0$ για κάθε $x \in \mathbb{R}$ και επειδή η $f$ είναι και συνεχής, είναι και γνησίως φθίνουσα.
Για τη συνάρτηση $f$, έχουμε ότι:
\begin{itemize}
    \item είναι συνεχής στο $[1,2]$ ως άθροισμα των συνεχών $y=-x+1$ (πολυωνυμική), $y=e^{-x}$ (εκθετική),
    \item $f(1) = \frac{1}{e} > 0$, $f(2) = -1+\frac{1}{e^2} < -1+\frac{1}{2^2} < 0$, οπότε $f(1) \cdot f(2) < 0$.
\end{itemize}
Επομένως, η συνάρτηση $f$ πληροί τις προϋποθέσεις του θεωρήματος Bolzano στο διάστημα $[1,2]$, άρα υπάρχει ένα τουλάχιστον $\rho \in (1,2)$ ώστε $f(\rho)=0$.
Επειδή η $f$ είναι γνησίως φθίνουσα έπεται ότι ο αριθμός $\rho$ είναι η μοναδική ρίζα της εξίσωσης $f(x)=0$ στο $\mathbb{R}$. Ακόμη είναι $\rho \in (1,2)$ άρα $\rho > 1$.

\item Το εμβαδό του ζητούμενου χωρίου $\Omega$ ισούται με $E(\Omega) = \int_{1}^{\rho} |f(x)|dx$.
Για κάθε $x \in (-\infty, \rho)$ είναι $x < \rho \xrightarrow{\text{f γνησ. φθίνουσα}} f(x) > f(\rho) \Rightarrow f(x) > 0$, άρα
\[ E(\Omega) = \int_{1}^{\rho} f(x)dx = \int_{1}^{\rho} \left(-x+1+e^{-x}\right)dx = \left[-\frac{(x-1)^2}{2}-e^{-x}\right]_{1}^{\rho} = -\frac{(\rho-1)^2}{2}-e^{-\rho}+e^{-1} \text{ τ.μ.} \]
Όμως ισχύει $f(\rho)=0 \Leftrightarrow -\rho+1+\frac{1}{e^\rho}=0 \Leftrightarrow e^{-\rho}=\rho-1$, οπότε
\[ E(\Omega) = -\frac{(\rho-1)^2}{2}-(\rho-1)+e^{-1} \text{ τετραγωνικές μονάδες}. \]
\end{alist}